\documentclass{report}
\usepackage{amsmath,amssymb}
\setlength{\parindent}{0mm}
\setlength{\parskip}{1em}
\begin{document}
\begin{center}
\rule{6in}{1pt} \
{\large Remi Lam \\
{\bf An Optimization Framework for Hybrid First-Principles Data-Driven Modeling}}

1539 Cambridge st \\ Cambridge \\ MA 02139 \\ USA
\\
{\tt rlam@mit.edu}\\
Lior Horesh\\
Haim Avron\\
Karen Willcox\end{center}

Mathematical models are used extensively for diverse tasks including
analysis, optimization, and decision making.
Frequently, those models are principled but imperfect representations of reality.
This is either due to incomplete physical description of the underlying
phenomenon (simplified governing equations, defective boundary
conditions, etc.), or due to numerical approximations (discretization,
linearization, round-off error, etc.).
Model misspecification can lead to erroneous model predictions, and
respectively suboptimal decisions associated with the intended end-goal
task.
To mitigate this effect, one can amend the available model using limited
data produced by experiments or higher fidelity models.
A large body of research has focused on estimating explicit model parameters.
This work takes a different perspective and targets the construction of a
correction model operator with implicit attributes.
We investigate the case where the end-goal is inversion and illustrate
how appropriate choices of properties imposed upon the correction and
corrected operator lead to improved end-goal insights.


\end{document}
